\subsection{Построение модели AS-IS}\label{subsec:business-process-as-is-modeling}

С точки зрения организации-исполнителя автоматизируемый бизнес-процесс разработки отдельного веб-сервиса в модели AS-IS включает в себя следующие этапы:

\begin{enumerate}
	\item Заказчик находит дефект или неточность в работе веб-сервиса и сообщает об этом руководителю проекта.
	\item Руководитель проекта анализирует полученную от заказчика обратную связь и создаёт задачу в трекере.
	\item Руководитель команды знакомится с поставленной задачей в трекере, определяет сложность и риски. Если выполнение задачи подразумевает существенные ограничения, то заказчик информируется о невозможности реализации.
	\item Руководитель команды назначает ответственных за выполнение задачи.
	\item Поставленный разработчик вносит изменения в исходный код веб-сервиса и отправляет результат во временную ветку удалённого репозитория.
	\item Руководитель команды производит ручное тестирование веб-сервиса, проверяет качество кода в ветке удалённого репозитория и даёт обратную связь. Если по задаче возникают замечания, то она отправляется на доработку.
	\item Руководитель команды сливает исходный код в основную ветку удалённого репозитория.
	\item Системный администратор производит ручную сборку исходного кода веб-сервиса и запускает полученный артефакт в staging-среде.
	\item Руководитель проекта тестирует изменения на staging-сервере. Если по задаче возникают замечания, то она отправляется на доработку.
	\item Руководитель проекта демонстрирует изменения заказчику. Если по задаче возникают замечания, то она отправляется на доработку.
	\item Системный администратор производит ручную сборку исходного кода веб-сервиса и запускает полученный артефакт в production-среде.
	\item Заказчик подтверждает руководителю проекта тот факт, что дефект или неточность в работе веб-сервиса была исправлена.
\end{enumerate}
